\subsubsection*{Training Procedure}

\subsubsection*{Modeling Objective}

The three-regime framework was developed to evaluate whether regime-specific modeling can improve predictive performance relative to a single global model. This design was motivated by residual diagnostics from the base model, which indicated non-uniform error structure across the target space.

The target variable remains surface soil moisture at 5 cm depth, denoted as \texttt{soil\_moisture\_5cm}.

\subsubsection*{Regime Definition}

Regimes were defined using thresholds derived from the empirical distribution of the target variable.

\begin{equation}
t_1 = 0.20, \qquad t_2 = 0.313
\end{equation}

\begin{table}[H]
\centering
\begin{tabular}{lll}
\toprule
\textbf{Regime} & \textbf{Condition} & \textbf{Description} \\
\midrule
Dry & $y \le t_1$ & Low-moisture, stable dynamics \\
Transition & $t_1 < y \le t_2$ & Mixed persistence and variability \\
Wet & $y > t_2$ & High-moisture, rapid dynamics \\
\bottomrule
\end{tabular}
\caption{Regime definitions based on empirical quantiles.}
\end{table}

These thresholds serve as the primary regime boundaries throughout the analysis.

\subsubsection*{Model Architecture}

The framework follows a hard-gated mixture-of-experts \cite{jacobs1991moe} architecture.

\begin{table}[H]
\centering
\begin{tabular}{ll}
\toprule
\textbf{Component} & \textbf{Role} \\
\midrule
Anchor model & Trained on full dataset; provides global prediction \\
Transition specialist & Trained on transition-regime samples \\
Wet specialist & Trained on wet-regime samples \\
Dry regime & Handled directly by anchor model \\
\bottomrule
\end{tabular}
\caption{Components of the three-regime architecture.}
\end{table}

A dedicated dry specialist was not retained in the final design. Empirically, the anchor model sufficiently captured dry-regime behavior, and introducing an additional model did not justify the added complexity.

\subsubsection*{Feature Usage}

The anchor model is trained on a shared base feature set. Transition and wet specialists use augmented inputs that include both regime-specific features and the anchor prediction:

\begin{equation}
x_{\text{specialist}} = \big[x_{\text{regime}}, \hat{y}_{\text{anchor}}\big]
\end{equation}

This formulation allows specialist models to learn corrections relative to a global estimate rather than reconstructing the full signal.

\subsubsection*{Training Configurations}

\begin{table}[H]
\centering
\begin{tabular}{p{4cm} p{9cm}}
\toprule
\textbf{Configuration} & \textbf{Description} \\
\midrule
Oracle-gated & Uses ground-truth regimes for routing; provides upper-bound performance \\
Self-gated (double-pass) & Uses anchor predictions for routing; deployable setting \\
Residual double-pass & Specialists predict residual corrections instead of absolute values \\
\bottomrule
\end{tabular}
\caption{Training configurations evaluated for the three-regime framework.}
\end{table}

\subsubsection*{Oracle-Gated Training}

In the oracle-gated setting, regime membership is determined using ground-truth values. The anchor model is trained on the full dataset, while transition and wet specialists are trained on their respective regime subsets.

During evaluation, samples are routed using true regime labels. Although not deployable, this configuration provides an estimate of the maximum achievable benefit under perfect gating.

\subsubsection*{Self-Gated Training}

In the self-gated setting, regime labels are derived from model predictions rather than ground truth. To reduce leakage, pseudo-labels are generated using out-of-fold predictions from a time-series cross-validation scheme.

At inference time, the anchor model produces an initial prediction, which is used to assign regime membership. The corresponding specialist model is then applied.

\subsubsection*{Residual Specialist Training}

In the residual variant, specialist models predict corrections to the anchor prediction:

\begin{equation}
r = y - \hat{y}_{\text{anchor}}, \qquad
\hat{y} = \hat{y}_{\text{anchor}} + \hat{r}
\end{equation}

This formulation simplifies the learning task by focusing each specialist on regime-specific residual patterns.

\subsubsection*{Temporal Structure}

All experiments use a time-based split to preserve temporal ordering:

\begin{table}[H]
\centering
\begin{tabular}{ll}
\toprule
\textbf{Split} & \textbf{Time Range} \\
\midrule
Training & Jan 2017 -- Dec 2020 \\
Validation & Jan 2021 -- Dec 2022 \\
Test & Jan 2023 -- Dec 2025 \\
\bottomrule
\end{tabular}
\caption{Temporal split used for three-regime experiments.}
\end{table}

All training and gating procedures respect this temporal structure to prevent leakage.

\subsubsection*{Interpretation}

The three-regime framework serves both as a predictive model and as an analytical tool. Its primary purpose is to test whether explicitly separating dry, transition, and wet dynamics can address the regime-dependent error structure observed in the base model.

While this decomposition aligns with the underlying physical intuition, its effectiveness ultimately depends on reliable regime identification and sufficient data within each regime.
